\documentclass{article}
\usepackage[utf8]{inputenc}
\usepackage{amsmath}
\usepackage{geometry}

\title{İntegral Çözüm Raporu}
\author{}
\date{\today}

\begin{document}

\maketitle

\section*{Soru}
Aşağıdaki integralin $D_n$ katsayısı için hesaplanması:
\[
D_n = \frac{2}{l} \int_0^l x \sin\left(\frac{n\pi}{l}x\right) \, dx
\]

\section*{Çözüm}
Bu integral \textbf{Kısmi İntegrasyon} yöntemi ile çözülür. Formülümüz: $\int u \, dv = uv - \int v \, du$.

\subsection*{Değişken Seçimi}
\begin{align*}
u &= x \implies du = dx \\
dv &= \sin\left(\frac{n\pi}{l}x\right) \, dx \implies v = -\frac{l}{n\pi} \cos\left(\frac{n\pi}{l}x\right)
\end{align*}

\subsection*{İşlemler}
İntegrali formülde yerine koyalım:
\[
\int_0^l x \sin\left(\frac{n\pi}{l}x\right) \, dx = \left[ -x \frac{l}{n\pi} \cos\left(\frac{n\pi}{l}x\right) \right]_0^l + \frac{l}{n\pi} \int_0^l \cos\left(\frac{n\pi}{l}x\right) \, dx
\]

Sınırları yerine koyarsak:
\begin{itemize}
    \item Birinci kısım: $-\frac{l^2}{n\pi}\cos(n\pi) - 0$
    \item İkinci kısım: $\sin(n\pi) = 0$ olduğundan sonuç 0'dır.
\end{itemize}

\section*{Sonuç}
Başlangıçtaki $\frac{2}{l}$ katsayısını ekleyerek sonucu buluruz:
\[
D_n = \frac{2}{l} \left( -\frac{l^2}{n\pi} \cos(n\pi) \right) = -\frac{2l}{n\pi} (-1)^n
\]
Veya işareti düzenlersek:
\[
\boxed{D_n = \frac{2l}{n\pi} (-1)^{n+1}}
\]

\end{document}